\section{Introduction}\label{intro}

Bid rigging is among the most costly forms of anti-competitive behavior in public procurement: estimates from exposed cases place the average cartel overcharge between 15 and 30 percent of contract value, and the OECD conservatively attributes to collusion an inefficiency wedge on the order of one fifth of public purchases in affected markets \citep{oecd2015, bosio2022public}. Because explicit evidence of coordination is almost always absent from the administrative record, a growing literature has sought to infer collusion indirectly from the statistical properties of bidding data. Classical detection screens exploit the fact that cartel agreements compress the distribution of bids: colluding firms submit prearranged prices that cluster tightly around the designated winner, yielding low bid dispersion, few active bidders, and predictable rotation patterns \citep{porter1993detecting, pesendorfer2000study, bajari2003detecting, conley2016detecting, imhof2018screening}. A complementary and influential approach, recently advanced by \citet{kawai2023using}, focuses on the threshold between winning and losing: under competition, the observable characteristics of narrowly winning and narrowly losing firms should be continuous at a bid margin of zero, and any statistically sharp discontinuity at the cutoff identifies an anti-competitive wedge. Applied to Japanese construction procurement, this Regression Discontinuity (RD) design uncovers sharp compression in bidder backlog and strong incumbency at the cutoff---signatures consistent with the canonical model of rotation-based collusion.

A first-order question for this literature is whether its screens travel. The cumulative evidence supporting canonical screen predictions is drawn almost exclusively from cartels operating in high-income economies with deep formal labor markets, dense bidder pools, and strong institutional capacity for antitrust enforcement. Whether the same statistical signatures appear in cartels operating in thinner institutional environments is an open question---and one with direct policy stakes, because emerging-market competition authorities increasingly rely on off-the-shelf screening tools developed and validated abroad to triage investigations \citep{oecd2015, bosio2022public}. If canonical screens systematically mislabel cartels under these conditions, existing enforcement pipelines may be misfiring in precisely the settings where deterrence capacity is weakest.

This paper leverages a rare opportunity to answer that question directly. Between 2009 and 2019, Brazil's federal antitrust authority---the Conselho Administrativo de Defesa Econ\^{o}mica (CADE)---prosecuted a series of procurement cartels operating in the state of S\~{a}o Paulo, securing administrative convictions against 35 firms across seven distinct cartel cases, with case files identifying the specific product markets and time windows over which each ring was active. We match these convicted firms, and the tender windows and product markets delineated in CADE's administrative records, to the universe of 3.8~million individual bids placed on S\~{a}o Paulo's electronic procurement platform, the \emph{Bolsa Eletr\^{o}nica de Compras} (BEC), over the same period. We further link the resulting bid-level dataset to the universe of formal-sector employment relationships in Brazil through the Rela\c{c}\~{a}o Anual de Informa\c{c}\~{o}es Sociais (RAIS), enabling us to characterize the workforce composition and worker mobility histories of every firm in the BEC at the time of each bid. To our knowledge, this is the first cartel-detection study to combine auction-level behavioral data, employer--employee matched labor records, and formal antitrust convictions as a single integrated benchmark.

Our research design exploits two complementary benchmark samples. The first is a \emph{pair-matched} sample of 107 pregão auctions in which CADE rulings directly implicate both the winning and the runner-up firms as co-conspirators in the same cartel case during the same active window---the cleanest ground truth available for any large-scale procurement-cartel study to date. The second is a \emph{cartel-active} sample of roughly 2{,}350 auctions in which at least one convicted firm participated during the active period of its cartel. The pair-matched sample delivers precision at the cost of statistical power; the cartel-active sample delivers power at the cost of a noisier treatment. Reporting results in parallel across both benchmarks disciplines our inference: a claim that survives both lenses is credible, a claim that appears in only one demands interpretation. We refer to the two samples collectively as \emph{conviction-anchored detection benchmarks}.

Our first set of results documents a striking failure of canonical screens in this institutional setting. On the pair-matched sample of 107 auctions, classical variance screens reverse their expected sign: relative to the population of comparable clean auctions, bid spread, the coefficient of variation of bids, and the number of active bidders are all \emph{higher}---not lower---at cartel-tainted auctions, with two-sided $p$-values below 0.02 in every case. These inversions are the opposite of the compression predicted by the canonical bid-rotation model that motivates classical screens and underpins \cite{kawai2023using}. On the broader cartel-active sample, where power is not a binding constraint, the same classical screens recover their canonical sign but lose virtually all discriminative power: the areas under the receiver-operating-characteristic curves (AUCs) for bid spread, coefficient of variation, and bidder count all lie between 0.46 and 0.48, statistically indistinguishable from chance. A direct replication of the Kawai close-margin RD on the cartel-active subsample yields no statistically detectable discontinuity in incumbency or last-bid behavior at the cutoff, and a density test in the spirit of \citet{imhof2018screening} detects no manipulation of the running variable. In short, neither the canonical variance-based screens nor the RD screen of \cite{kawai2023using} detect the cartels that CADE has formally convicted.

Our second set of results shows that network-based screens constructed from administrative labor data recover strong and statistically robust cartel detection on the same benchmarks. For each auction, we construct two network screens from the five-year window of RAIS records preceding the tender: the number of workers who have moved between the winner and the runner-up firms, and the Jaccard similarity of the two firms' historical workforces. These yield AUCs of 0.825 and 0.815 respectively on the cartel-active benchmark, outperforming every classical screen by wide margins. A third screen, built from the co-bidding network of shared tender participations, yields an AUC of 0.69, providing independent corroboration through a distinct information channel. Critically, the worker-flow signal is neither an artifact of sectoral composition nor a mechanical re-expression of co-bidding intensity. Restricting attention to firm pairs that share the same 2-digit sectoral classification leaves the worker-flow AUC essentially unchanged at 0.8249, and the residual of worker-flow after partialling out the co-bidding density of each pair still achieves an AUC of 0.74. The conviction-anchored benchmark does not merely declare cartels detectable in S\~{a}o Paulo; it pinpoints the channel through which the detection signal propagates---worker mobility, not bid dispersion.

We interpret these findings through the lens of institutional thinness. In settings where the universe of firms capable of bidding on a given product is small and the formal-sector labor market for the specialized workers that serve those bids is shallow, coordination among colluding firms generates naturally detectable traces in worker flows: personnel rotations, key-employee transfers, and shared specialized labor become vehicles for, and therefore signatures of, the collusive communication that substitutes for the statistical-pattern coordination captured by canonical screens in thicker markets. This mechanism---which we develop formally in a simple model in Section~\ref{coll}---rationalizes both the failure of bid-dispersion screens and the success of network-based screens in the Brazilian setting, and it predicts that classical screen-based enforcement strategies will systematically under-detect cartels in precisely the institutional environments where their deterrence value would be highest.

This paper makes three contributions. \emph{First}, methodologically, we provide the first stress test of the RD cartel-detection framework of \citet{kawai2023using} against an independent criminal ground truth, and we document that its canonical predictions can reverse sign in institutional environments with thin bidder pools. This is not a repudiation of the close-margin RD design, which remains internally valid; rather, it is a warning that the \emph{direction} of its test statistic depends on primitives that are not invariant across settings. \emph{Second}, substantively, we introduce conviction-anchored detection benchmarks by linking CADE rulings, BEC bid-level records, and RAIS employer--employee data, and we show that worker-flow and co-bidding network screens constructed from these sources substantially outperform classical screens in detecting convicted cartels. Because such data are increasingly available to national competition authorities in middle-income economies, the screens we propose are directly deployable at negligible marginal cost. \emph{Third}, conceptually, we advance the view that collusion leaves institution-specific signatures, and that the design of an effective detection screen should depend on which coordination technology is cheapest for cartel members to operate in the relevant environment---a framing that organizes and partially reconciles the growing patchwork of contradictory empirical findings in the detection literature \citep{conley2016detecting, imhof2018screening, kawai2023using}.

The remainder of the paper is structured as follows. Section~\ref{hist} describes the institutional setting of S\~{a}o Paulo's public procurement and the CADE cartel cases that anchor our ground truth. Section~\ref{literature} situates our contribution in the procurement-collusion detection literature and in the adjacent literatures on worker mobility and law-and-economics enforcement. Section~\ref{data} describes the BEC, RAIS, and CADE data sources and the construction of the pair-matched and cartel-active benchmark samples. Section~\ref{coll} presents our empirical strategy---the RD replication, the portfolio of classical screens, and the network-based screens---and sketches a simple theoretical framework linking screen predictions to institutional primitives. Section~\ref{res} reports the main results. Section~\ref{concl} discusses policy implications and concludes.

% TODO-BIB: add to biblio.bib before final compile:
%  - cattaneo2018manipulation (Cattaneo, Jansson, Ma 2018/2020 density test) [used in §\ref{coll}, §\ref{res}]
%  - huber2019imhof (Huber & Imhof 2019, JCLE — machine learning screens)
%  - chassang2022ortner (Chassang & Ortner 2022/2023 — missing trader tests)
%  - clark2018houde (Clark, Coviello, Gauthier, Shneyerov 2018 — retail gasoline collusion)
%  - jager2023heining (Jäger, Heining & Kline — worker flow literature)
%  - morselli2009 (Morselli — covert networks)
%  - wachs2019network (Wachs, Kertesz — network-based cartel detection)
%  - baltrunaite2020 (Baltrunaite et al. — political connections in procurement)
%  - harrington2008 (Harrington — optimal cartel enforcement)
%  - miller2009 (Miller — strategic leniency)
